
\section{Benchmark Details} \label{app:benchmarks}

Table~\ref{tab:app-benchmarks} lists the eight test sets of xRouteBench, their categories, sizes, and metrics, totaling 4{,}767 test queries. The Generic LLM Tasks track further decomposes into 13 subtasks, and the memory datasets are retrieved with top-5 RAG over turn pairs.

\begin{table}[h]
\centering
\caption{\textbf{The eight test sets of xRouteBench.} Sizes are the number of test queries; metrics are exact match (EM), multiple-choice accuracy (MC), token-level F1, execution-based code pass rate, math answer matching, and a persona-conditioned LLM judge.}
\label{tab:app-benchmarks}
\small
\resizebox{\linewidth}{!}{%
\begin{tabular}{lllrl}
\toprule
\textbf{Category} & \textbf{Test set} & \textbf{Content} & \textbf{\#Test} & \textbf{Metric} \\
\midrule
Generic LLM Tasks & Generic mix & 13 subtasks & 3{,}729 & EM/MC/F1/GSM8K/MATH/code \\
\midrule
\multirow{2}{*}{Memory} & LoCoMo & long-conversation QA & 314 & F1 \\
 & LongMemEval & long-term memory QA & 101 & F1 \\
\midrule
TimeSeries & TimeSeries & 7 reasoning skills & 127 & MC \\
\midrule
\multirow{3}{*}{Vision} & Geometry3K & geometry math (image) & 61 & EM \\
 & MathVista & visual math reasoning & 100 & EM/MC \\
 & Charades-Ego & egocentric video & 27 & EM \\
\midrule
Personalized & Chatbot Arena / MT-Bench & preference prompts & 308 & LLM judge \\
\midrule
\textbf{Total} & & & \textbf{4{,}767} & \\
\bottomrule
\end{tabular}
}
\end{table}

Table~\ref{tab:app-classic} details the 13 subtasks that make up the Generic LLM Tasks track.

\begin{table}[h]
\centering
\caption{\textbf{Composition of the Generic LLM Tasks track.} The table lists the 13 subtasks, their target skills, and the number of test queries, totaling 3{,}729 examples.}
\label{tab:app-classic}
\small
\setlength{\tabcolsep}{6pt}
\begin{tabular}{llr}
\toprule
\textbf{Subtask} & \textbf{Skill} & \textbf{\#Test} \\
\midrule
MBPP & code generation & 500 \\
MATH & mathematical reasoning & 500 \\
GSM8K & mathematical reasoning & 500 \\
MMLU-Pro & knowledge QA & 500 \\
OpenBookQA & knowledge QA & 500 \\
ARC-Challenge & knowledge QA & 500 \\
MMLU & knowledge QA & 500 \\
CommonsenseQA & commonsense QA & 50 \\
BoolQ & commonsense QA & 50 \\
SQuAD & reading comprehension & 50 \\
HellaSwag & commonsense QA & 50 \\
HumanEval & code generation & 16 \\
AIME (2020--2024) & competition math & 13 \\
\bottomrule
\end{tabular}
\end{table}

\subsection{Generic LLM Tasks} \label{app:general_nlp}

The Generic LLM Tasks track is deliberately a mixture rather than a single task family. It places knowledge questions, commonsense inference, reading comprehension, mathematical reasoning, and code generation behind the same routing interface. All examples use text input, but an ARC-Challenge question, an AIME problem, and a Python synthesis task reward very different capabilities and impose different output constraints. The track therefore asks whether a router can recognize these distinctions within an apparently uniform text setting, rather than defaulting to one model for every natural language request. Table~\ref{tab:app-classic} lists the 13 source benchmarks and their test sizes.

We serialize each example as one fixed text query before collecting candidate responses. A system instruction specific to the task states the required response format, and the user problem follows it. Choice items include their options. Mathematics items request a final answer in the expected form. Code generation items provide the programming task and tests, while SQuAD supplies its passage and question. We score each response with the source benchmark's native objective. We use choice accuracy for knowledge and commonsense tasks, token-level F1 for SQuAD, answer matching for GSM8K, MATH, and AIME, and execution-based pass rates for MBPP and HumanEval.

\subsection{Long Context Conversational Memory} \label{app:memory}

The Memory track fixes retrieval and asks a narrower routing question. Given the same retrieved conversational evidence, which model is most likely to use it correctly? LoCoMo and LongMemEval contain questions whose supporting facts may be separated from the query by many turns or sessions. Some can be answered by recovering one explicit fact, whereas others require combining facts across sessions or resolving a later update against an earlier statement. This distinction is useful for routing because it separates the cost of carrying a long history from the ability to interpret the evidence selected from it.

We construct every memory query through the same retrieval pipeline. Each conversation history is divided into adjacent turn pairs that retain speaker and date information. A fixed Contriever \citep{izacard2022unsupervised} encoder embeds the question and the turn pairs, and the five most similar pairs are inserted into a brief answer prompt. LongMemEval additionally includes the date of the question, which provides the temporal reference needed for its time-sensitive items. Every candidate therefore receives the same retrieved evidence for a query. Differences in score reflect its use of that evidence rather than a different retrieval result. We evaluate both datasets with token-level F1.

\subsection{Time Series Pattern Reasoning} \label{app:timeseries}

The TimeSeries track is built from TSRBench and covers anomaly detection, similarity analysis, noise understanding, pattern recognition, inductive reasoning, causality analysis, and event prediction. These problems are numerical, but they often require recognizing structure that is easier to see as a shape than as a list of values. A local spike may signal an anomaly, while periodicity, changes in trend, or agreement between two series emerge over a longer range. The track asks whether a router can distinguish the models that handle these different forms of temporal reasoning well.

We convert every series into a common text query before collecting candidate responses. First, we render each series as a line chart and pass the chart to a fixed Gemma-3-27B-IT \citep{team2025gemma} captioner, which produces a paragraph description of its visible temporal pattern. We append these descriptions and the raw numerical values to the original question and its answer options. The raw values preserve exact numerical evidence and are truncated after 200 values when necessary, while the descriptions make higher-level structure explicit. The charts are used only in this offline captioning step. Every candidate model receives the same text query rather than an image. Models return an option letter, and we score the responses with multiple-choice accuracy.







\subsection{Visual Mathematical Reasoning}  \label{app:math_image}

The Visual Mathematical Reasoning track routes image-grounded mathematics drawn from Geometry3K \citep{lu2021inter} and MathVista \citep{lu2024mathvista}, where a geometry diagram or a scientific figure carries information the question depends on. The strongest candidate on these problems can differ from the strongest on text mathematics, so the track probes whether a router follows per-model strength as the query type changes. We render each problem to text before it reaches the pool, describing its image with a fixed vision-language model and appending the description to the question, so every candidate reasons over one identical query. Rendering the image once holds perception constant across the pool and turns visual mathematics into a query the whole pool can answer, so the track scores the reasoning quality of each candidate on the same input. These rendered queries read as machine-written descriptions of a figure, a distribution that departs from the natural questions the routers are trained on, so the track also tests how well a router generalizes when the query type shifts.

A single frozen vision-language model, Gemma-3-27B-IT by default, produces every description, so all candidates see the same rendering of a given image. Its prompt asks it to report every visible number, symbol, angle, length, and relationship and to withhold the solution, which keeps the answer out of the query and leaves the reasoning to the routed model. The returned text is appended to the original problem, from which Geometry3K's image placeholder is removed, to form the routed query. Geometry3K is graded by math answer matching against its numeric answer, and MathVista is graded by exact match on its open-ended items and by multiple-choice accuracy on its choice items, following the question-type field of each query. Every query keeps its ground-truth answer and, where present, its answer choices, so the scored responses populate the same query--model matrix of quality and cost that supervises every other track (\S\ref{sec:protocol}). Figure~\ref{fig:multimodal_examples} shows one example from each dataset with the description Gemma-3-27B-IT produces for it.

\begin{figure}[H]
\centering
\begin{minipage}[c]{0.24\linewidth}
\centering
\includegraphics[width=0.9\linewidth]{figs/geometry3k.png}\\[3pt]
{\footnotesize (a) Geometry3K, which asks for $x$.}
\end{minipage}
\hfill
\begin{minipage}[c]{0.72\linewidth}
\begin{mybox}[Gemma-3-27B-IT description]
\footnotesize The diagram shows a circle with two intersecting chords. One chord is divided into segments of length $4$ and $6$, and the other into segments of length $x$ and $8$. The label $x$ marks one of these segments, and the two chords cross at a point inside the circle.
\end{mybox}
\end{minipage}

\vspace{6pt}

\begin{minipage}[c]{0.40\linewidth}
\centering
\includegraphics[width=\linewidth]{figs/mathvista.png}\\[3pt]
{\footnotesize (b) MathVista, which asks for the spring compression $d$.}
\end{minipage}
\hfill
\begin{minipage}[c]{0.56\linewidth}
\begin{mybox}[Gemma-3-27B-IT description]
\footnotesize A block of mass $m$ slides to the left along a horizontal frictionless surface with velocity $v$ and approaches a coiled spring of constant $k$ fixed to a wall. Labels mark where the block first touches the spring and where it stops, and $d$ denotes the distance between these two points. A caption states that the spring force does negative work, decreasing speed and kinetic energy.
\end{mybox}
\end{minipage}
\caption{\textbf{One example from each dataset in the Visual Reasoning track, shown with the description that Gemma-3-27B-IT produces for its image.} Geometry3K (a) provides a geometry diagram, and MathVista (b) provides a scientific figure. Each description is appended to the problem text to form the query the router sees.}
\label{fig:multimodal_examples}
\end{figure}


\subsection{Multi-View Video Recognition} \label{app:video}

The video track builds on Charades-Ego \citep{sigurdsson2018charades}, which records each activity simultaneously from a first-person egocentric camera worn by the actor and a third-person exocentric camera facing the scene (Figure~\ref{fig:charades_sample}). The two viewpoints carry different evidence, and a deployment often captures only one of them, so the routing decision here depends on which views a query provides. We describe every available view in text and route the merged description, so the whole candidate pool answers the same query whether one view or both are present. This makes the video track the regime in xRouteBench where the available modality varies from query to query, and it tests whether a router still selects the right model when the visual evidence is partial.

We build three classification tasks from the annotations, predicting the activity, the verb, or the object of the depicted action, each answered by a compact identifier drawn from the task label inventory, with the verb and object labels recovered from the Charades action mapping. For each paired clip we match the two views by their shared identifier and select the action segment whose egocentric and exocentric occurrences overlap most closely in time, which keeps both views on the same moment, and we then keep the first-person view, the third-person view, or both at random to populate the single-view and dual-view regimes. From each retained view we sample frames inside the aligned window and pass them to the vision-language model, which returns a structured description of the actor's motion, the handled objects, and the scene. These descriptions are merged into one query that states the task, lists the candidate identifiers with their names, and requests an identifier as the answer, holding every candidate to the same text input.

\begin{figure}[H]
\centering
\setlength{\tabcolsep}{2pt}
\renewcommand{\arraystretch}{0.7}
\begin{tabular}{c ccccc}
 & \multicolumn{5}{c}{\footnotesize time $\rightarrow$} \\
\rotatebox{90}{\footnotesize First-person} & \includegraphics[width=0.18\linewidth]{figs/charades/ego_t1.jpg} & \includegraphics[width=0.18\linewidth]{figs/charades/ego_t2.jpg} & \includegraphics[width=0.18\linewidth]{figs/charades/ego_t3.jpg} & \includegraphics[width=0.18\linewidth]{figs/charades/ego_t4.jpg} & \includegraphics[width=0.18\linewidth]{figs/charades/ego_t5.jpg} \\
\rotatebox{90}{\footnotesize Third-person} & \includegraphics[width=0.18\linewidth]{figs/charades/exo_t1.jpg} & \includegraphics[width=0.18\linewidth]{figs/charades/exo_t2.jpg} & \includegraphics[width=0.18\linewidth]{figs/charades/exo_t3.jpg} & \includegraphics[width=0.18\linewidth]{figs/charades/exo_t4.jpg} & \includegraphics[width=0.18\linewidth]{figs/charades/exo_t5.jpg} \\
\end{tabular}

\vspace{4pt}
\begin{mybox}[Gemma-3-27B-IT description]
\footnotesize \textit{Placeholder for the structured description the vision-language model produces from the available views.}
\end{mybox}
\caption{\textbf{A multi-view sample from Charades-Ego \citep{sigurdsson2018charades}, recorded from a first-person egocentric camera (top) and a third-person exocentric camera (bottom).} Each row is a viewpoint and each column is a sampled time step. Our transformation samples these frames, describes the available views with a vision-language model, and routes the resulting text query, randomly withholding one view to form single-view and dual-view queries.}
\label{fig:charades_sample}
\end{figure}

\subsection{Personalized Dialogue Preference} \label{app:personalized}

The Personalized track combines open-ended dialogue from MT-Bench and Chatbot Arena. Each query retains its original message sequence, including preceding system, user, and assistant messages. MT-Bench contributes instruction following conversations in which a later turn depends on an earlier exchange, while Chatbot Arena contributes natural user requests across diverse conversational settings. Rather than assuming that one answer is universally better, the track constructs supervision from the preference that a particular user profile would express after seeing two answers.

The user profiles used for preference elicitation are drawn from PersonaHub \citep{tencent2024personahub}. We sample 200 personas from this resource for data collection. Figure~\ref{fig:personahub_personas} presents ten examples from the full persona set used during data collection.

\begin{figure}[H]
\centering
\begin{minipage}{0.98\linewidth}
\begin{mybox}[Ten example personas used to condition the judge]
\footnotesize
\begin{tabular}{@{}p{0.47\linewidth}@{\hspace{0.03\linewidth}}p{0.47\linewidth}@{}}
\textbf{1.} A 71-year-old retired nurse from Italy, volunteering in hospice care and advocating for compassionate end-of-life support. & \textbf{6.} A 34-year-old scientist from London who is a social media influencer. \\
\textbf{2.} A 54-year-old divorced mother from Spain, running a successful winery and promoting sustainable viticulture practices. & \textbf{7.} A 41-year-old scientist from London who loves hiking. \\
\textbf{3.} A 63-year-old retired teacher from China, teaching calligraphy and preserving the art form for future generations. & \textbf{8.} An 87-year-old World War II veteran from Poland, sharing stories of his experiences and advocating for peace. \\
\textbf{4.} A 68-year-old retired engineer from Japan, practicing ikebana and teaching the art to younger generations. & \textbf{9.} A 31-year-old social worker from Colombia, supporting victims of domestic violence and fighting for gender equality. \\
\textbf{5.} A 21-year-old photographer from Paris who spends weekends volunteering. & \textbf{10.} A 23-year-old aspiring musician from Brazil, fusing traditional and modern sounds and promoting cultural exchange through music. \\
\end{tabular}
\end{mybox}
\end{minipage}
\caption{\textbf{Ten examples from the 200 PersonaHub personas sampled for preference collection.} For each comparison, DeepSeek-V3.1 receives one selected persona as its role specification when judging the two candidate answers. The candidate models do not receive the persona.}
\label{fig:personahub_personas}
\end{figure}

For each dialogue, we sample two models at random from the candidate pool of 18 models and generate one response from each under the same conversation history. DeepSeek-V3.1 is conditioned on the selected persona to act as the judge, comparing the two answers and returning a preference for either response or a tie. The persona is supplied to the judge only, not to either candidate model. Each comparison is converted into a pairwise supervision record whose label is the final judged user preference. The router learns from these preference outcomes rather than from the persona description itself. This construction keeps the candidate responses comparable while making the routing target sensitive to the user utility represented by the judge.



\section{Router Details} \label{app:routers}

The routers in \method share the candidate pool and task interface described above, but they expose different information to the routing decision. Rule-based baselines select from a fixed property of the pool, while single-turn methods make one model selection from the query and logged routing outcomes. Multi-turn methods additionally condition on intermediate calls and may decide whether further routing is useful. Personalized methods add a user and interaction history, so their target is the model a particular user is most likely to prefer rather than one model that is best on average. Table~\ref{tab:app-routers} summarizes the 17 built-in implementations along these differences.

\begin{table}[t]
\centering
\caption{\textbf{Built-in routers in \method, grouped by routing family.} For each method, the table summarizes the information available to the routing decision (State) and the rule used to select or aggregate candidate models (Selection).}
\label{tab:app-routers}
\small
\setlength{\tabcolsep}{4pt}
\renewcommand{\arraystretch}{1.08}
\newcommand{\routerspace}{\addlinespace[2pt]}
\begin{tabular}{
    p{0.20\linewidth}
    p{0.30\linewidth}
    p{0.32\linewidth}
}
\toprule
\textbf{Router} & \textbf{State} & \textbf{Selection} \\
\midrule
\rowcolor{barGray} \multicolumn{3}{c}{\textit{Rule-based baselines}} \\
Smallest-LLM
& candidate parameter counts
& always selects the smallest candidate \\
\routerspace
Largest-LLM
& candidate parameter counts
& always selects the largest candidate \\
\midrule
\rowcolor{barGray} \multicolumn{3}{c}{\textit{Single-turn routers}} \\
kNNRouter
& query embedding and nearby logged queries
& votes over the models preferred by nearest neighbors \\
\routerspace
SVMRouter
& query embedding
& kernel classifier predicts a candidate \\
\routerspace
MLPRouter
& query embedding
& MLP classifier predicts a candidate \\
\routerspace
MFRouter
& query and model latent factors
& ranks candidates by their interaction score \\
\routerspace
EloRouter
& logged pairwise model outcomes
& always selects the highest-rated candidate \\
\routerspace
RouterDC
& query and candidate representations
& contrastive query--model matching score \\
\routerspace
Hybrid LLM
& query embedding and a small/large model pair
& predicts whether the small model is sufficient \\
\routerspace
AutoMix
& small-model draft and verification signal
& accepts the draft or escalates to the large model \\
\routerspace
GraphRouter
& query--model interaction graph
& predicts performance on query--model edges \\
\routerspace
CausalLM Router
& textual query and candidate list
& generates the selected model name \\
\midrule
\rowcolor{barGray} \multicolumn{3}{c}{\textit{Multi-turn routers}} \\
Router-R1
& query and accumulated search results
& iteratively searches specialists or terminates and aggregates \\
\routerspace
kNN-MultiRound
& sub-queries and their embeddings
& routes each sub-query with kNN and aggregates the answers \\
\routerspace
LLM-MultiRound
& textual query, decomposition, and candidate list
& an LLM chooses routes for sub-queries and aggregates \\
\midrule
\rowcolor{barGray} \multicolumn{3}{c}{\textit{Personalized routers}} \\
GMTRouter
& user, session, query, model, and response interactions
& predicts user-conditioned model preference \\
\routerspace
PersonalizedRouter
& user features, task description, query, and model
& predicts preference for a user--query pair \\
\bottomrule
\end{tabular}
\end{table}


\subsection{Rule-Based Baselines}

Smallest-LLM and Largest-LLM provide fixed reference points for the learned routers. They ignore the query and always select the candidate with the smallest or largest declared parameter count, respectively. These rules make no attempt to identify per-query model strengths, but they expose the two simple deployment policies against which query-aware routing should be compared: consistently favoring the smallest available model or consistently favoring the largest one.

\subsection{Single-Turn Routers}

Single-turn routers terminate after one model selection, so their differences lie in how they represent a query and score candidates. EloRouter is query-independent, but it replaces a fixed parameter-count rule with a global ranking estimated from pairwise outcomes in the routing data. It therefore captures which model is strongest on average while deliberately discarding the variation between individual queries.

kNNRouter, SVMRouter, and MLPRouter instead make the selection query-specific from its embedding. kNNRouter retrieves similar training queries and transfers their observed best-model choices through a vote or distance-weighted vote, requiring no fitted decision function beyond the stored examples \citep{li2025rethinking,shnitzer2023large}. SVMRouter and MLPRouter fit discriminative boundaries over the same embedding space, using a kernel classifier and a multilayer perceptron, respectively. MFRouter takes a different view of the logged query--model matrix: it learns latent representations for both sides and selects the model with the strongest query--model interaction. These methods all learn from per-query outcomes, but differ in whether the candidate is represented by neighboring solved examples, a classifier label, or a learned latent factor.

The remaining single-turn methods enrich this compatibility score in different ways. RouterDC learns query--model matching with dual contrastive objectives over query--model, query--query, and cluster-level relations \citep{chen2024routerdc}. GraphRouter instead passes information over a query--model interaction graph and predicts the quality of an unobserved query--model edge \citep{feng2024graphrouter}. CausalLM Router verbalizes the query and the available candidates, then fine-tunes a causal language model to generate the selected model name \citep{ong2024routellm}. In this case, the query representation, candidate representation, and selection score are combined within the language model rather than implemented as separate embedding modules.

Hybrid LLM and AutoMix are cost-aware two-model cascades. Hybrid LLM learns from the quality gap between the smallest and largest candidates, then uses a thresholded prediction to decide whether the smaller model is sufficient \citep{ding2024hybrid}. AutoMix first obtains a draft from the smaller model and a verification signal for that draft; it retains the draft when the signal is reliable and otherwise escalates the query to the larger model \citep{aggarwal2024automix}. We list both methods with the single-turn family because they return one final model answer through a fixed cascade, rather than repeatedly choosing among arbitrary candidates and aggregating an open-ended history. Their extra calls are nevertheless part of the inference cost.

\subsection{Multi-Turn Routers}

Multi-turn routers treat intermediate responses as part of the routing state. Router-R1 is a pretrained routing agent that reasons over the query and the results gathered so far. At each step it can issue a \texttt{<search>} call to a suitable specialist, incorporate the returned evidence, or terminate and produce an aggregate answer; its policy is trained with trajectory-level reinforcement learning rather than pointwise best-model labels \citep{zhang2025router}. The cost of its reasoning, search, and final aggregation calls is included in the routed trajectory.

The two multi-round baselines use decomposition more explicitly. kNN-MultiRound first breaks a complex query into a small set of sub-queries, applies the kNN routing rule to each one, and combines the resulting answers. LLM-MultiRound uses an LLM to express both the decomposition and the model choice in text, then executes the selected sub-queries and aggregates their outputs. Both methods can assign different candidates to different parts of a problem, but they also introduce decomposition and aggregation calls that a one-shot router does not pay for.

\subsection{Personalized Routers}

Personalized routers replace the single global notion of quality with a user-conditioned preference. GMTRouter represents users, sessions, queries, candidate models, and responses as nodes in a heterogeneous interaction graph, allowing a decision for the current query to draw on earlier interactions from the same user \citep{xie2025gmtrouter}. PersonalizedRouter similarly uses a graph-based scorer, while explicitly incorporating user features and task descriptions alongside the query and candidate representations \citep{dai2025personalizedrouter}. Both methods are trained from comparisons between candidate answers, so a high score means that a model is predicted to be preferred for this user and query, not simply that it has the highest population-average task score.





\section{Candidate Pool and Pricing} \label{app:pricing}

Table~\ref{tab:app-pricing} lists the 18 candidate models and their per-token prices. Prices are in USD per 1M tokens, with input and output priced separately; the blended average is their mean. The input price spans a roughly $25\times$ range, from \$0.05 to \$1.25 per 1M tokens. The 17-model no-cogito pool drops the most expensive model (cogito-v2-1-671b).

\begin{table}[h]
\centering
\caption{\textbf{The 18 candidate LLMs, sorted by blended average price.} Prices in USD per 1M tokens.}
\label{tab:app-pricing}
\small
\setlength{\tabcolsep}{6pt}
\begin{tabular}{rlrrrl}
\toprule
\textbf{\#} & \textbf{Model} & \textbf{Params} & \textbf{Input} & \textbf{Output} & \textbf{Service} \\
\midrule
1 & gemma-2-9b-it & 9B & 0.10 & 0.10 & NVIDIA \\
2 & llama-3-8b-instruct-lite & 8B & 0.10 & 0.10 & Together \\
3 & gpt-oss-20b & 20B & 0.05 & 0.20 & Together \\
4 & rnj-1-instruct & 15B & 0.15 & 0.15 & Together \\
5 & mistral-7b-instruct-v0.3 & 7B & 0.20 & 0.20 & NVIDIA \\
6 & mistral-small-3-24b-instruct & 24B & 0.10 & 0.30 & Together \\
7 & qwen2.5-7b-instruct & 7B & 0.20 & 0.20 & NVIDIA \\
8 & qwen2.5-7b-instruct-turbo & 7B & 0.30 & 0.30 & Together \\
9 & gpt-oss-120b & 120B & 0.15 & 0.60 & Together \\
10 & llama-4-maverick & 402B & 0.27 & 0.85 & Together \\
11 & mixtral-8x7b-instruct-v0.1 & 46.7B & 0.60 & 0.60 & NVIDIA \\
12 & qwen3-next-80b-a3b-instruct & 80B & 0.15 & 1.50 & Together \\
13 & qwen3-coder-next & 200B & 0.50 & 1.20 & Together \\
14 & llama-3.3-70b-instruct-turbo & 70B & 0.88 & 0.88 & Together \\
15 & llama3-70b-instruct & 70B & 0.90 & 0.90 & NVIDIA \\
16 & deepseek-v3.1 & 671B & 0.60 & 1.70 & Together \\
17 & mixtral-8x22b-instruct-v0.1 & 140.6B & 1.20 & 1.20 & NVIDIA \\
18 & cogito-v2-1-671b & 671B & 1.25 & 1.25 & Together \\
\bottomrule
\end{tabular}
\end{table}

\section{Human Preference Collection on Slack} \label{app:slack}

The human preference dataset of \S\ref{sec:extension-study} is collected through a Slack application built on the OpenClaw server of \method, as illustrated in Figure~\ref{fig:slack}. A user asks a question directly in Slack, and the system samples two LLMs at random from the ten-candidate pool and generates one answer with each. The conversation view then presents the two responses as anonymized Answer A and Answer B, with their positions randomly shuffled, and the user clicks a button to mark A better, B better, or a tie. In multi-turn sessions the user follows up freely, and every turn is labeled in the same way. Fifteen users contributed 40 sessions and 234 pairwise preference records in total.
\begin{figure}[h]
\centering
\includegraphics[width=1.0\linewidth]{figs/slack_platform.png}
\caption{\textbf{Slack interface for collecting pairwise user preferences.} Two anonymized model responses are shown as Answer A and Answer B in randomized order, and users select A, B, or a tie for each interaction turn.}
\label{fig:slack}
\end{figure}

\section{Multi-Agent Topologies} \label{app:mas}

Table~\ref{tab:app-mas} summarizes the five coordination topologies used in \S\ref{sec:extension-study}. Star, Tree, Graph, and Chain follow MultiAgentBench \citep{zhu2025multiagentbench}, and Plan-Exec-Sum follows the static router design of GraphPlanner \citep{feng2026graphplanner} with width 3 and depth 1. Star and Plan-Exec-Sum differ in two respects: in Star, the planner produces free-form subtasks and consolidates the actors' outputs itself, whereas in Plan-Exec-Sum the planner decomposes the query into three atomic sub-queries and a dedicated summarizer merges the executors' answers. Every node is replaced by a router that receives the node's prompt and selects a model from the candidate pool for that call. We adopt the topology of GraphPlanner without training its planner, since the training requires live answers from candidate models that have since been retired.

\begin{table}[h]
\centering
\caption{\textbf{The five multi-agent topologies, their structures, and the number of LLM calls per query.}}
\label{tab:app-mas}
% \small
\resizebox{\linewidth}{!}{%
\begin{tabular}{llc}
\toprule
\textbf{Topology} & \textbf{Structure} & \textbf{LLM calls} \\
\midrule
Star & planner decomposes $\rightarrow$ 3 actors in parallel $\rightarrow$ planner consolidates & 6 \\
Tree & root planner $\rightarrow$ 2 sub-planners refine $\rightarrow$ 2 actors $\rightarrow$ root consolidates & 7 \\
Graph & 3 actors answer independently $\rightarrow$ one full-communication revision round & 7 \\
Chain & 3 agents relay sequentially, each verifying and improving the previous answer & 4 \\
Plan-Exec-Sum & planner emits 3 atomic sub-queries $\rightarrow$ 3 executors $\rightarrow$ summarizer merges & 6 \\
\bottomrule
\end{tabular}
}
\end{table}

\section{Prompt Usage} \label{app:prompt}

This section reports the fixed prompt templates used to construct and
evaluate xRouteBench, as well as the templates used by routers that make
language-model calls internally. Curly brackets denote instance-dependent
fields. Each prompt is presented in a titled box, combining the fixed
instruction and the instance fields in the order in which the model reads
them, rather than splitting them into separate system and user blocks. In the
implementation, the fixed instruction is sent as a system message when the
API supports that role; otherwise, the two parts are concatenated with
\texttt{[System Instruction]} and \texttt{[User Query]} delimiters. Thus, all
candidate models receive the same logical prompt for a given query.

\newcommand{\prompttext}[1]{{\ttfamily\noindent\detokenize{#1}\par}}
\newcommand{\promptfield}[1]{{\ttfamily\detokenize{{#1}}}}

\subsection{Benchmark-Query Templates}

\paragraph{Generic LLM Tasks.}
The 13 Generic LLM Tasks use seven task-specific templates. The
multiple-choice instruction is shared by MMLU, MMLU-Pro, BoolQ, and
HellaSwag; CommonsenseQA, OpenBookQA, and ARC-Challenge use the same
instruction but retain their source choice formatting.

\begin{mybox}[Standard multiple choice]
\footnotesize
\prompttext{Answer the following multiple-choice question by selecting the correct option
(A, B, C, or D). You MUST put your final answer letter in a parenthesis.}

\medskip

\prompttext{Question:
{question}}

\prompttext{Options:
A. {choice_1}
B. {choice_2}
C. {choice_3}
D. {choice_4}}
\end{mybox}

\begin{mybox}[Source-formatted multiple choice]
\footnotesize
\prompttext{Answer the following multiple-choice question by selecting the correct option
(A, B, C, or D). You MUST put your final answer letter in a parenthesis.}

\medskip

\prompttext{Question: {question}
(A) {choice_1}
(B) {choice_2}
(C) {choice_3}
...}
\end{mybox}

\begin{mybox}[GSM8K]
\footnotesize
\prompttext{Answer the following math question step by step.}
\medskip
\prompttext{Question: {question}}
\end{mybox}

\begin{mybox}[MATH and AIME]
\footnotesize
\prompttext{Answer the following math question. Make sure to put the answer (and only
answer) inside \boxed{}.}
\medskip
\prompttext{Question: {question}}
\end{mybox}

\begin{mybox}[MBPP]
\footnotesize
\prompttext{You are an expert Python programmer. Implement the function with no irrelevant
words or comments. Put your code in this format: [BEGIN] <Your Code> [Done]}
\medskip
\prompttext{Task: {task_description}}

\prompttext{Your code should pass these tests: {tests}}

\end{mybox}

\begin{mybox}[HumanEval]
\footnotesize
\prompttext{You are an expert Python programmer. Implement the function body only. Do not
repeat the function signature or docstring. Put the code in this format:
[BEGIN] <Your Code - function body only> [Done]}
\medskip
\prompttext{Complete the following function: {function_signature_and_docstring}}


\end{mybox}

\begin{mybox}[SQuAD]
\footnotesize
\prompttext{Answer the question using the provided context. If the answer is not in the
context, output noanswer.}

\prompttext{{source query containing the context and question}}
\end{mybox}

\paragraph{Conversational memory.}
For both memory datasets, Contriever retrieves the top five turn-pair chunks
before the following prompt is formed. The retrieval result is fixed across
candidate models for each query. LoCoMo uses:
\begin{mybox}[LoCoMo]
\footnotesize
\prompttext{You answer questions about conversation history. Reply with a short phrase
only.}

\prompttext{Below is relevant information from the conversation history:}

\prompttext{{retrieved turn-pair chunks}}
\medskip
\prompttext{Based on the above context, write an answer in the form of a short phrase for
the following question. Answer with exact words from the context whenever
possible.}

\prompttext{Question: {question} Short answer:}
\end{mybox}
If no chunk is retrieved, the context field is replaced by \texttt{No
relevant information available.} LongMemEval uses the analogous template,
with its question date made explicit:
\begin{mybox}[LongMemEval]
\footnotesize
\prompttext{You answer questions based on chat history. Reply with a short phrase only.}

\prompttext{I will give you several history chats between you and a user. Please answer
the question based on the relevant chat history.}
\medskip
\prompttext{History Chats:}

\prompttext{### Session {index}:
Session Content:
{retrieved chunk}}

\prompttext{...}

\prompttext{Current Date: {current_date}}
\prompttext{Question: {question}}
\prompttext{Short Answer:}
\end{mybox}
Here, an empty history is represented by \texttt{No relevant chat history
available.}

\paragraph{Visual mathematical reasoning.}
Geometry3K and MathVista first use one frozen vision-language model to turn
the image into text. The captioning call is:
\begin{mybox}[Image captioning]
\footnotesize
\prompttext{Describe the image for solving the problem.
Include all visible text, numbers, symbols, angles, lengths, and geometric
relationships.
Be concise and factual. Do not include introductory phrases.}

\prompttext{{image}}
\end{mybox}
The resulting caption is then given to every candidate model with the
following answer prompt:
\begin{mybox}[Visual mathematical reasoning]
\footnotesize
\prompttext{You are an expert in solving math and geometry problems. Think step-by-step,
and when you are ready to provide the final answer, you MUST enclose it in
\boxed{}. }
\medskip
\prompttext{For example, if the answer is 42, write \boxed{42}. If it is a
multiple choice question and the option is C, write \boxed{C}.}
\medskip
\prompttext{Image Description: {caption}}

\prompttext{Answer the following question based on the provided image description:
{question}}
\end{mybox}

\paragraph{Time-series reasoning.}
We first render every series as an image and request a one-paragraph caption.
The captioner receives the following system instruction and one instruction
sampled uniformly from 20 equivalent phrasings (the first phrasing is shown):
\begin{mybox}[Time-series captioning]
\footnotesize
\prompttext{You are a time series captioner.}

\prompttext{Write a paragraph that analyzes the time series.}

\prompttext{{time-series plot}}
\end{mybox}
The other phrasings differ only in wording (e.g., ``Create a detailed
description of the time series in one paragraph'') and impose the same
one-paragraph description requirement. The answer prompt combines the caption
with the raw values:
\begin{mybox}[Time-series reasoning]
\footnotesize
\prompttext{You are a time series analysis expert. Answer the multiple-choice question by
providing the correct option letter (e.g., A, B, C, or D). Be concise.}
\medskip
\prompttext{1. Time Series Description: (1) {series_name}: {caption}, ...}
\prompttext{2. Raw time series: (1) {series_name}: [{values}], ...}
\medskip
\prompttext{based on the one of the following information to reason the final answer:
{original_question}}
\end{mybox}
For perception and causality-analysis items, this two-source block replaces
the \texttt{<ts><ts/>} marker in the source question; it is appended for the
other two task types.

\paragraph{Multi-view video recognition.}
Five frames from each available egocentric or exocentric view are summarized
independently before classification. The video captioner receives:
\begin{mybox}[Video captioning]
\footnotesize
\prompttext{You are given multiple frames sampled from a short time window of a video.}
\medskip
\prompttext{Return ONLY a JSON object with the following keys (no extra text):}
\prompttext{  "motion": string,}
\prompttext{  "objects": [string],}
\prompttext{  "summary": string}
\medskip
\prompttext{Definitions:}
\medskip
\prompttext{motion: describe what the person is doing (hands/arms/body) in one short phrase.}
\prompttext{objects: list the main target objects the person is acting on / interacting with (not background clutter).}
\prompttext{summary: a short description of the scene in one sentence.}



\prompttext{{video frames}}
\end{mybox}
Each candidate model then receives the two structured captions, the relevant
label inventory, and this classification prompt:
\begin{mybox}[Video classification]
\footnotesize
\prompttext{You are a helpful AI assistant tasked with answering questions about videos
based on descriptions. }
\prompttext{Read the provided text carefully and return ONLY the
ID (e.g. c001, v003, o012) representing the correct answer.}
\medskip
\prompttext{Task: Identify the {activity/verb/object} from video descriptions.}
\prompttext{Determine the correct id using the structured descriptions below.}
\medskip
\prompttext{First-person view description: {ego_caption}}
\prompttext{Third-person view description: {exo_caption}}
\prompttext{Available {activity/verb/object}s:
{id}: {label}
...}
\medskip
\prompttext{Provide your answer as the ID (e.g., c001, v003, o012).}
\end{mybox}

\paragraph{Personalized dialogue preference.}
For MT-Bench, the original multi-turn user/assistant message sequence is
preserved; for Chatbot Arena, the original user message is preserved. No
additional fixed instruction is inserted before a candidate answers. Two
candidate responses are then compared by a judge conditioned on a sampled
persona:
\begin{mybox}[Persona preference judge]
\footnotesize
\prompttext{You are simulating the following user persona:
{persona}.}

\prompttext{As this persona, evaluate which assistant response you would personally
prefer.}
\medskip
\prompttext{[User Question]
{question}}

\prompttext{[Assistant 1]
{answer_1}}

\prompttext{[Assistant 2]
{answer_2}}
\medskip
\prompttext{Output ONLY one of the following tokens exactly:
1
2
Tie
Do not output anything else.}
\end{mybox}

\subsection{Router-Internal Templates}

The preceding templates construct the query--model response matrix used by
all routers. Most single-turn routers consume this matrix without another LLM
prompt. The following templates are used only by the routers that perform
additional language-model calls during routing or aggregation.

\paragraph{kNN-MultiRound and LLM-MultiRound.}
Both multi-round routers decompose a query, obtain sub-answers, and aggregate
them. kNN-MultiRound selects a model for each sub-query with nearest-neighbour
lookup, whereas LLM-MultiRound asks an LLM to choose the model. Their shared
decomposition template is:
\begin{mybox}[Multi-round decomposition]
\footnotesize
\prompttext{Given the query '{query}', decompose it into as many as 4 meaningful
sub-queries (minimum 1, maximum 4).}

\prompttext{Try to cover the full scope of the original query by breaking it down into
multiple specific and distinct sub-tasks whenever possible.}

\medskip
\prompttext{Aim for the maximum number of high-quality sub-queries without introducing
redundancy.}

\prompttext{Each sub-query should be clear, self-contained, and semantically coherent.}
\medskip
\prompttext{Output only the decomposed sub-queries, one per line.
Do not include any other text, explanations, or headers in your output.
Each line should contain only one sub-query.}
\end{mybox}
LLM-MultiRound appends the candidate names and descriptions, then replaces the
last three lines by the following routing constraint:
\begin{mybox}[Multi-round model selection]
\footnotesize

\prompttext{You will then be provided with descriptions of the following Large Language
Models (LLMs): {model_list}.}
\medskip
\prompttext{{model_descriptions}}
\medskip
\prompttext{For each sub-query, select the single LLM that is most likely to generate the
highest-quality response, regardless of cost or efficiency.}

\prompttext{Focus entirely on maximizing effectiveness and providing the most accurate
and relevant output. Base your decision strictly on the descriptions of the models.}
\medskip
\prompttext{Output only the decomposed sub-queries and the full name of the selected LLM
for each. Output exactly one sub-query and one LLM per line.}
\medskip
\prompttext{Each line must be formatted as follows:}

\prompttext{<sub-query>: <LLM name>}
\medskip
\prompttext{Use a colon ':' as the separator. Do not include any other text, explanations, or headers in your output.}

\end{mybox}
Each selected model is queried with:
\begin{mybox}[Multi-round specialist]
\footnotesize

\prompttext{You are a helpful assistant. You are participating in a multi-agent reasoning process, where a base model
delegates sub-questions to specialized models like you.}
\medskip
\prompttext{Your task is to do your absolute best to either:}

\prompttext{Answer the question directly, if possible, and provide a brief explanation;}

\prompttext{Offer helpful and relevant context, background knowledge, or insights related
to the question, even if you cannot fully answer it.}
\medskip
\prompttext{If you are completely unable to answer the question or provide any relevant or
helpful information, you must:}

\prompttext{Clearly state that you are unable to assist with this question.}

\prompttext{Explicitly instruct the base model to consult other LLMs for further
assistance.}
\prompttext{Keep your response clear, concise, and informative (preferably under 512
tokens).}

\prompttext{Stay strictly on-topic.}

\prompttext{Do not include irrelevant or generic content.}
\medskip
\prompttext{Here is the sub-question for you to assist with: {sub_query}}

\end{mybox}
For non-multiple-choice tasks, the aggregator receives:
\begin{mybox}[Multi-round aggregation]
\footnotesize

\prompttext{You are given a question along with auxiliary information, which consists of
several sub-questions derived from the original question and their respective
answers.}

\prompttext{Use this information to answer the original question if relevant, but make
your own reasoning step by step before arriving at the final answer.}
\medskip
\prompttext{Important: Your final answer MUST be clearly marked and enclosed within
<answer> and </answer> tags at the end of your response. No other part of the output should be inside these tags.}
\medskip
\prompttext{Auxiliary Information: {sub_query_and_response_pairs}}

\prompttext{Question: {query}}

\prompttext{Let's think step by step.}

\end{mybox}
For multiple-choice tasks, it instead receives:
\begin{mybox}[Multi-round multiple-choice aggregation]
\footnotesize

\prompttext{You are given a multiple-choice question and supporting sub-answers. Use the information only if helpful.}
\medskip
\prompttext{Question: {original_query}}

\prompttext{Supporting information: {sub_query_and_response_pairs}}
\medskip
\prompttext{Rules:}

\prompttext{Select exactly one option: A, B, C, D, or E.}

\prompttext{Output only the letter in <answer> tags.}

\prompttext{No explanation.}
\medskip
\prompttext{Format:}

\prompttext{<answer>A</answer>}

\end{mybox}

\paragraph{Router-R1.}
Router-R1 appends descriptions of the available models to the following
template. Its agent can request a specialist model within \texttt{<search>}
tags; the returned text is inserted as \texttt{<information>} in the next
round.
\begin{mybox}[Router-R1]
\footnotesize

\prompttext{Answer the given question.}

\prompttext{Every time you receive new information, you must first conduct reasoning
inside <think> ... </think>.}

\prompttext{After reasoning, if you find you lack some knowledge, you can call a
specialized LLM by writing a query inside
<search> LLM-Name:Your-Query </search>.}
\medskip
\prompttext{STRICT FORMAT RULES for <search>:}

\prompttext{Use an exact model name from: {candidate_model_names}.}

\prompttext{Replace Your-Query with a concrete question.}

\prompttext{Never copy model descriptions into <search>.}

\prompttext{Never output the placeholder format literally.}

\prompttext{Before each LLM call, reason about why external information is needed and
which model is appropriate.}

\prompttext{The returned response appears between <information> and </information>.}

\prompttext{Different models may be called multiple times.}
\medskip
\prompttext{{model_descriptions}}
\medskip
\prompttext{If no further external knowledge is needed, output the final answer inside
<answer> ... </answer>.}

\prompttext{Do not output empty answer tags.}
\medskip
\prompttext{Question: {question}}

\end{mybox}

\paragraph{CausalLM.}
The CausalLM router is trained to complete the selected model name after this
prefix:
\begin{mybox}[CausalLM]
\footnotesize
\prompttext{You are an intelligent router that selects the best Large Language Model (LLM)
for a given query.}
\medskip
\prompttext{Available LLMs: {model_1}, {model_2}, ...}

\prompttext{Based on the query content, complexity, and requirements, predict which LLM
would provide the best response.}
\medskip
\prompttext{Query: {query}}

\prompttext{Best LLM:}
\end{mybox}

\paragraph{AutoMix.}
AutoMix first asks a small model to answer with the following template:
\begin{mybox}[AutoMix answer generation]
\footnotesize
\prompttext{You are given a question. Answer the question as concisely as you can, using a
single phrase if possible.}
\medskip
\prompttext{Question: {question}}

\prompttext{Answer: The answer is'}
\end{mybox}
It uses the following few-shot verifier:
\begin{mybox}[AutoMix verifier]
\footnotesize
\prompttext{Question: Whose lost work was discovered in a dusty attic in 1980?}

\prompttext{AI Generated Answer: Shakespeare}

\prompttext{Instruction: Your task is to evaluate if the AI Generated Answer is correct,
based on the provided question. Provide the judgement and reasoning for each
case. Choose between Correct or Incorrect.}

\prompttext{Evaluation: The lost work of Shakespeare was discovered in 1980 in a dusty
attic.}

\prompttext{Verification Decision: The AI generated answer is Correct.}

\prompttext{---}

\prompttext{Question: In which month does the celestial event, the Pink Moon, occur?}

\prompttext{AI Generated Answer: July}

\prompttext{Instruction: Your task is to evaluate if the AI Generated Answer is correct,
based on the provided question. Provide the judgement and reasoning for each
case. Choose between Correct or Incorrect.}

\prompttext{Evaluation: The Pink Moon is unique to the month of April.}

\prompttext{Verification Decision: The AI generated answer is Incorrect.}

\prompttext{---}

\prompttext{Question: Who is believed to have painted the Mona Lisa in the early 16th
century?}

\prompttext{AI Generated Answer: Vincent van Gogh}

\prompttext{Instruction: Your task is to evaluate if the AI Generated Answer is correct,
based on the provided question. Provide the judgement and reasoning for each
case. Choose between Correct or Incorrect.}

\prompttext{Evaluation: The Mona Lisa was painted by Leonardo da Vinci in the early 16th
century.}

\prompttext{Verification Decision: The AI generated answer is Incorrect.}

\prompttext{---}

\prompttext{Question: How far away is the planet Kepler-442b?}

\prompttext{AI Generated Answer: 1,100 light-years}

\prompttext{Instruction: Your task is to evaluate if the AI Generated Answer is correct,
based on the provided question. Provide the judgement and reasoning for each
case. Choose between Correct or Incorrect.}

\prompttext{Evaluation: The Kepler-442b is located 1,100 light-years away.}

\prompttext{Verification Decision: The AI generated answer is Correct.}

\prompttext{---}

\prompttext{Question: {question}}

\prompttext{AI Generated Answer: {generated_answer}}

\prompttext{Instruction: Your task is to evaluate if the AI Generated Answer is correct,
based on the provided question. Provide the judgement and reasoning for each
case. Choose between Correct or Incorrect.}

\prompttext{Evaluation:}
\end{mybox}
The verifier's estimated correctness determines whether the query is escalated
to the larger model.

\paragraph{OpenClaw deployment router.}
When the deployed server uses prompt-based routing, it selects one model with:
\begin{mybox}[OpenClaw deployment router]
\footnotesize

\prompttext{You are an intelligent LLM router. Choose the most suitable model for the user's query.}
\medskip
\prompttext{Available models: {model_descriptions}}
\medskip
\prompttext{Rules:}

\prompttext{1. Simple greetings/daily chat -> cheaper models (8b, 9b size)}

\prompttext{2. Q&A/knowledge retrieval -> chatqa models}

\prompttext{3. Instruction following/structured output -> mistral models}

\prompttext{4. Code generation/technical questions -> nemotron or larger models}

\prompttext{5. Complex reasoning/deep analysis -> 70b or larger models}
\medskip
\prompttext{IMPORTANT: Only return the model name, nothing else!}
\medskip
\prompttext{Model names: {model_list}}

\prompttext{User query: {query}}

\end{mybox}
For image and video inputs, OpenClaw first uses the following preprocessing
prompts; the returned text is included in the router query.

\begin{mybox}[OpenClaw image and video preprocessing]
\footnotesize
\textbf{Image input:}

\prompttext{Describe this image concisely in 2--3 sentences.}

\textbf{Video input:}

\prompttext{Describe what you see in these video frames.}
\end{mybox}
\subsection{Multi-Agent Topology Templates}

The multi-agent topologies in \S\ref{app:mas} route every LLM-call node
independently from the prompt assigned to that node. Repeated nodes use the
same template and differ only in their indexed fields. The final node is shared
by all five topologies and is given after the topology-specific templates.

\paragraph{Star.}
The central planner decomposes the problem for \promptfield{n_actors} actors,
then consolidates their reports.

\begin{mybox}[Star planner]
\footnotesize
\prompttext{You are the central planner of a team of {n_actors} agents working on the problem
below. Decompose the problem into exactly {n_actors} subtasks, one per agent, so
that their combined results solve the problem. Output EXACTLY {n_actors} lines,
one subtask per line, no numbering or extra words.}
\medskip
\prompttext{Problem:
{user_query}}
\end{mybox}

\begin{mybox}[Star actor]
\footnotesize
\prompttext{You are agent {i+1} in a team. Solve the following subtask assigned by your
planner. Be concise and factual.}
\medskip
\prompttext{Original problem:
{user_query}}

\prompttext{Your subtask:
{subtask}}
\end{mybox}

\begin{mybox}[Star planner consolidation]
\footnotesize
\prompttext{You are the central planner. Your agents reported the results below. Consolidate
them into a single coherent analysis of the problem.}
\medskip
\prompttext{Problem:
{user_query}}

\prompttext{Agent reports:
Subtask: {subtask_1}
Result: {result_1}}

\prompttext{Subtask: {subtask_2}
Result: {result_2}}

\prompttext{...}
\end{mybox}

\paragraph{Tree.}
The root planner assigns two complementary sub-problems to team leads. Each
lead refines its branch into one task for a worker, after which the root planner
consolidates the two reports.

\begin{mybox}[Tree root planner]
\footnotesize
\prompttext{You are the top-level planner. Split the problem below into exactly 2
complementary sub-problems for your two subordinate team leads. Output EXACTLY
2 lines, one sub-problem per line, no numbering.}
\medskip
\prompttext{Problem:
{user_query}}
\end{mybox}

\begin{mybox}[Tree sub-planner]
\footnotesize
\prompttext{You are team lead {b+1}. Refine the sub-problem below into one concrete task
for your worker agent. Output ONE line only.}
\medskip
\prompttext{Original problem:
{user_query}}

\prompttext{Your sub-problem:
{branch_task}}
\end{mybox}

\begin{mybox}[Tree worker]
\footnotesize
\prompttext{You are a worker agent. Solve the task below concisely and factually.}
\medskip
\prompttext{Original problem:
{user_query}}

\prompttext{Your task:
{task}}
\end{mybox}

\begin{mybox}[Tree root consolidation]
\footnotesize
\prompttext{You are the top-level planner. Consolidate your two teams' reports into a
single coherent analysis of the problem.}
\medskip
\prompttext{Problem:
{user_query}}

\prompttext{Team reports:
Sub-problem: {branch_task_1}
Task: {task_1}
Result: {result_1}}

\prompttext{Sub-problem: {branch_task_2}
Task: {task_2}
Result: {result_2}}
\end{mybox}

\paragraph{Graph.}
Each of \promptfield{n_actors} agents first answers independently. In each
subsequent communication round, an agent receives the other agents' current
answers and revises its own answer.

\begin{mybox}[Graph actor: initial round]
\footnotesize
\prompttext{You are agent {i+1} of {n_actors} independent agents. Solve the problem below
on your own. Show brief reasoning, then your answer.}
\medskip
\prompttext{Problem:
{user_query}}
\end{mybox}

\begin{mybox}[Graph actor: revision round]
\footnotesize
\prompttext{You are agent {i+1} of {n_actors}. You communicated with all other agents;
their current answers are below. Reconsider and give your revised reasoning
and answer.}
\medskip
\prompttext{Problem:
{user_query}}

\prompttext{Your previous answer:
{answers[i]}}

\prompttext{Agent {j+1}'s current answer:
{answers[j]}}

\prompttext{...}
\end{mybox}

\paragraph{Chain.}
The first agent answers independently. Each later agent verifies and improves
the immediately preceding decision before passing its result onward.

\begin{mybox}[Chain: first agent]
\footnotesize
\prompttext{You are agent 1 in a chain of {n_agents} agents. Solve the problem below. Show
brief reasoning, then your answer. Your output will be passed to the next agent.}
\medskip
\prompttext{Problem:
{user_query}}
\end{mybox}

\begin{mybox}[Chain: later agent]
\footnotesize
\prompttext{You are agent {i+1} in a chain of {n_agents} agents. The previous agent passed
you their decision below. Verify it, fix any mistakes, and pass on your
improved reasoning and answer.}
\medskip
\prompttext{Problem:
{user_query}}

\prompttext{Previous agent's decision:
{decision}}
\end{mybox}

\paragraph{Plan-Exec-Sum.}
The GraphPlanner-style topology first decomposes the query into
\promptfield{width} atomic sub-queries. Each executor receives its sub-query
without an additional instruction, and the summarizer combines the results.

\begin{mybox}[Plan-Exec-Sum planner]
\footnotesize

\prompttext{You are a query decomposition assistant. Your task is to decompose the user's query into exactly {width} atomic and
independent sub-queries.}
\medskip
\prompttext{Keep each sub-query self-contained and non-overlapping.}

\prompttext{Prefer factual, directly answerable units.}

\prompttext{Output EXACTLY {width} lines, one sub-query per line, no numbering or extra
words.}

\medskip

\prompttext{User query: {user_query}}

\end{mybox}

\begin{mybox}[Plan-Exec-Sum executor]
\footnotesize
\prompttext{{sub_query}}
\end{mybox}

\begin{mybox}[Plan-Exec-Sum summarizer]
\footnotesize
\prompttext{You are a professional summarizer.
Summarize the following content into a concise, coherent paragraph without
bullet points.
Make it fluent and logically connected.}
\medskip
\prompttext{Content:
Sub-query: {sub_query_1}
Answer: {answer_1}}

\prompttext{Sub-query: {sub_query_2}
Answer: {answer_2}}

\prompttext{...}

\prompttext{Summary:}
\end{mybox}

\paragraph{Shared final node.}
Every topology uses the same final node. Its system instruction is restored
from the original task's \texttt{[System Instruction]} block, so the required
answer format remains task-specific. When there is team context, the final
node receives the following user prompt:

\begin{mybox}[Shared final node: user prompt with team context]
\footnotesize
\prompttext{{user_query}}

\prompttext{Below is supporting analysis from your team. Use it if helpful, but follow the
answer format required above.}

\prompttext{[Team Analysis]
{context}}
\end{mybox}

Without team context, the final node receives \promptfield{user_query} alone.









\section{The ComfyUI Visual Interface} \label{app:comfyui}

\method exposes its full routing pipeline as a graph on the ComfyUI canvas, where every node is a library component and every edge is an artifact that flows between components. Two input nodes supply the source benchmarks and the candidate pool $\mathcal{M}$, a data-engine node consumes them and emits the query--model matrix as a single edge, and each router node consumes that matrix and emits its evaluation. The matrix records the per-query performance and cost of every candidate, the supervision that every router in \S\ref{sec:general-router} is trained on. The graph traces the evaluation protocol of \S\ref{sec:protocol} from left to right, and the router nodes appear in menu groups named after the three router families, so the taxonomy of \S\ref{sec:general-router} is visible in the node menu. The canvas therefore renders the architecture of Figure~\ref{fig:library} as a diagram a user reads and edits, and replacing a router replaces one node while the data-engine and evaluation nodes stay in place, so the modularity of the library becomes visible on the canvas.

The data nodes realize the three stages of the data engine. The Select Datasets and Select LLMs nodes declare the query source and the candidate pool, and the Generate Data node runs query curation, response collection, and scoring in one call and writes the query--model matrix to a self-contained data directory. The router nodes cover every built-in router and are grouped in the menu by router family, and each router node reads its defaults from the same YAML configuration that the command line uses and renders every hyperparameter as a typed widget whose value, range, and options come from that file, so a canvas node and its scripted counterpart run one configuration. On execution a router node writes a runtime configuration, dispatches training and evaluation through the shared router registry, and returns a summary of the query count, the success count, the average performance, and the routing distribution over candidates. The Generate Data node hashes the selected datasets, candidate pool, and sample size into a metadata record and reuses an existing data directory when the record and its files are unchanged, which skips the costly response-collection stage on repeated runs.

\begin{figure}[H]
\centering
\includegraphics[width=\linewidth]{figs/comfyui.png}
\caption{\textbf{The ComfyUI interface of \method.} The source benchmarks and the candidate pool enter at the left, the data-engine node produces the query--model matrix, and each router node consumes the matrix and reports its evaluation, so the graph traces the routing pipeline from data construction to evaluation. Each router node exposes the hyperparameters of its library configuration as typed widgets.}
\label{fig:comfyui}
\end{figure}